% META: {"id": "1NSI-RESEAUX-RE-C4", "chapitre": "1NSI-RESEAUX", "type_objet": "remediation", "capacites": ["P-ARCH-04A"], "capacites_codes": ["C4"], "mode_creation": "ex_nihilo", "sources_inspiration": [], "status": "needs_review"}

%---------------------------------------------------------------
% FICHE DE REMEDIATION — C4 : un capteur mesure, un actionneur agit
%---------------------------------------------------------------

\textbf{Ce qui bloque.} Tu intervertis les deux rôles, ou tu prêtes au capteur le pouvoir de
déclencher quelque chose. Un capteur ne décide rien et n'agit sur rien : il produit une
valeur. Ce qui déclenche, c'est le programme ; ce qui agit, c'est l'actionneur.

\textbf{À retenir.} Deux questions suffisent. « Qu'est-ce qui produit la valeur que le
programme lit ? » désigne le capteur. « Qu'est-ce qui change dans le monde quand le
programme a décidé ? » désigne l'actionneur. Entre les deux, il n'y a que du code.

\begin{exercice}{1NSI-RES-RE-C4-EX1}{1}{10}
Une serre connectée comporte : une sonde d'humidité du sol, une électrovanne d'arrosage, un
thermomètre, un moteur d'ouverture du toit, un capteur de luminosité, une lampe horticole, un
bouton poussoir « arrosage manuel » et un écran d'affichage.
\begin{enumerate}
  \item Classer chaque composant en capteur ou actionneur, en justifiant chaque choix par
        « il mesure » ou « il agit ».
  \item Un élève écrit : « le thermomètre déclenche l'ouverture du toit quand il fait trop
        chaud ». Réécrire la phrase en distinguant ce qui mesure, ce qui décide et ce qui
        agit.
  \item Le bouton poussoir et l'écran ne mesurent aucune grandeur physique. À quel rôle les
        rattacher malgré tout, et pourquoi ?
\end{enumerate}
\end{exercice}
